% sumber: doc/TRD_M1.md, deploy/, doc/P1_STORAGE_DESIGN.md, doc/P2_LEASE_DESIGN.md, go.mod, services/web-frontend/package.json dan nginx.conf
\chapter{Daftar Teknologi Terpilih}

\section{Bahasa dan \textit{Runtime}}

Seluruh \textit{service} backend ditulis dalam Go 1.26. Setiap \textit{service} adalah satu \textit{module} Go dengan Dockerfile sendiri. Tahap \textit{build} memakai \texttt{golang:1.26-alpine}, dan hasilnya dijalankan di atas \texttt{alpine:3.22}. Go dipilih karena lebih mudah dibaca dibandingkan Rust dan ekosistemnya sudah matang, termasuk pustaka gRPC, klien RabbitMQ, dan OpenTelemetry yang dipakai sistem ini.

\section{Komunikasi Antar-\textit{Service}}

Komunikasi internal memakai gRPC. gRPC dipilih karena latensinya lebih rendah daripada REST dan kontrak antar-\textit{service} ditetapkan lewat berkas \textit{protobuf}. REST hanya dipakai pada batas luar, yaitu antara \textit{client} dan Client-Facing API, serta antara \textit{Aggregator} dan sumber data BMKG dan PVMBG, karena kedua batas itu memang menuntut fleksibilitas REST.

Pembaruan \textit{real-time} ke browser memakai SSE, singkatan dari \textit{Server-Sent Events}, bukan WebSocket. Arah datanya satu arah dari server ke browser, dan browser menyambung ulang otomatis saat koneksi putus, sehingga fitur dua arah WebSocket tidak dibutuhkan.

\section{Penyimpanan}

Data disimpan di PostgreSQL 16 (\texttt{postgres:16-alpine}). Tabel \texttt{hazard\_events} memiliki kolom bertipe untuk sepuluh \textit{field} kanonik dan kolom JSONB \texttt{attributes} untuk atribut spesifik sumber, sehingga \textit{field} baru dapat tersimpan tanpa migrasi skema. Pada Milestone 2, satu pod PostgreSQL berjalan di tiap node dan data dibagi ke 16 shard dengan 2 salinan per shard. Pembagian dan penyalinan dilakukan sendiri oleh \textit{Aggregator} dengan library standar Go, tanpa produk basis data terdistribusi baru dan tanpa replikasi bawaan Postgres. Alasan dan alternatif yang ditolak ada di \texttt{doc/P1\_STORAGE\_DESIGN.md} dan Bab XI. \textit{Lease}, kursor \textit{polling}, dan token PVMBG disimpan di satu PostgreSQL terpisah, \textit{coordination-db}, yang dipilih karena memakai teknologi yang sudah ada dan waktu server basis data dapat menjadi sumber waktu (\texttt{doc/P2\_LEASE\_DESIGN.md}, bagian Alternatif dan keputusan, dan Bab XII). User DB memakai PostgreSQL yang terpisah. Perbandingan dengan relasional kolom tetap, \textit{document store}, dan \textit{key-value store} dibahas pada Bab IX.

\section{\textit{Message Broker}}

Distribusi \textit{event} memakai RabbitMQ 3.13 dengan \textit{fanout exchange} \texttt{hazard.events}. Setiap \textit{consumer} memiliki \textit{durable queue} sendiri yang terikat ke \textit{exchange} tersebut. Alasan memilih RabbitMQ dibandingkan Kafka dan NATS dibahas pada Bab X.

\section{Autentikasi}

Autentikasi \textit{client} memakai JWT dengan algoritma HS256 yang disimpan sebagai \textit{cookie} \textit{HttpOnly}, sehingga tidak dapat dibaca JavaScript. Akses ke PVMBG memakai OAuth 2.0 dengan rotasi \textit{refresh token}, sedangkan akses ke BMKG memakai \textit{API key}. Alasan tiap mekanisme dibahas pada Bab VIII.

\section{Frontend dan \textit{Web Server}}

Antarmuka dibangun dengan SvelteKit 3 dan Svelte 5 karena lebih sederhana daripada menulis HTML, CSS, dan JavaScript native secara manual. Hasil \textit{build} berupa berkas statis yang disajikan oleh nginx 1.27.

nginx juga meneruskan \textit{path} \texttt{/api} ke Client-Facing API. Dengan cara ini browser memanggil API pada \textit{origin} yang sama dengan halaman, sehingga \textit{cookie} autentikasi tidak membutuhkan konfigurasi CORS (konfigurasi ada di \texttt{services/web-frontend/nginx.conf}).

\section{Deployment dan \textit{Observability}}

Seluruh komponen dijalankan di Kubernetes distribusi k3s pada tiga VM VirtualBox, sesuai ketentuan Milestone 2 yang mewajibkan Kubernetes dengan minimal dua node berkernel terpisah. k3s dipilih karena disarankan spesifikasi dan ringan untuk VM berukuran 2 GiB (\texttt{doc/TOPOLOGY.md}). Manifest dikelola dengan Kustomize (\texttt{deploy/kustomization.yaml}) dan \textit{image} disimpan di \textit{registry} lokal di node-a. Docker Compose masih dipakai untuk pengembangan lokal. Setiap \textit{service} mengirim \textit{trace} lewat OpenTelemetry ke Jaeger 2.22, sehingga satu permintaan dapat dilacak melewati beberapa \textit{service}. Prometheus 3.6 menyediakan metrik laju permintaan, galat, dan latensi per \textit{service} yang ditampilkan pada tab Monitor Jaeger. Grafana 12.2 menampilkan metrik semua \textit{service} pada satu \textit{dashboard}, dan kube-state-metrics 2.17 menyediakan metrik pod (\texttt{deploy/observability/}). Seluruh versi ini tertera pada \texttt{image} di manifest tiap komponen.

\section{Karakteristik Arsitektur yang Diprioritaskan}

Pilihan teknologi di atas dan rancangan pada Bab VI sampai XV mengikuti karakteristik arsitektur yang diprioritaskan kelompok. Karakteristik diurutkan dari spesifikasi tiap problem, dari batasan yang tercatat di tiap bab, dan dari isi repositori. Gambar \ref{fig:karakteristik} menghubungkan karakteristik itu dengan keputusan arsitektur, gaya arsitektur, dan komponen logis.

\begin{figure}[H]
\centering
\includegraphics[width=\linewidth,height=0.5\textheight,keepaspectratio]{figures/karakteristik.pdf}
\caption{Karakteristik arsitektur, keputusan, gaya, dan komponen logis}
\label{fig:karakteristik}
\end{figure}

\begin{longtable}[]{|>{\raggedright\arraybackslash}p{0.16\linewidth}|>{\raggedright\arraybackslash}p{0.12\linewidth}|>{\raggedright\arraybackslash}p{0.5\linewidth}|>{\raggedright\arraybackslash}p{0.1\linewidth}|}
\hline
Karakteristik & Prioritas & Wujud pada sistem & Bab \\
\hline
\endhead
\textit{Availability} & Utama & Jalur \textit{ingest} dipisah dari jalur akses. Sumber yang mati dijawab dengan data terakhir dan \texttt{meta.stale\_since}, dan 503 hanya bila \textit{store} kosong. & VII \\
\hline
\textit{Fault tolerance} & Utama & Poller per sumber dengan satu permintaan berjalan, \textit{circuit breaker}, dan \textit{health check} yang melaporkan \texttt{degraded}. Jaeger bukan dependensi, sehingga \textit{service} tetap melayani saat Jaeger mati. & VII \\
\hline
\textit{Performance} & Utama & gRPC antar-\textit{service}, goroutine dengan I/O non-blocking, \textit{load shedding} pada batas \textit{inflight}, dan \textit{rate limit} per \textit{client}. Uji p95 dan \textit{load testing} ada di \texttt{tests/load}. & VII \\
\hline
\textit{Security} & Utama & Kredensial terpisah per instansi, JWT dalam \textit{cookie} \textit{HttpOnly}, \textit{scope} yang menyembunyikan field Mentah dari Media, dan \textit{network} internal untuk tiap penyimpanan. & VIII, IX \\
\hline
\textit{Agility} & Pendukung & Kolom JSONB menyerap atribut baru tanpa migrasi, dan \textit{consumer} baru cukup mengikat \textit{queue} ke \textit{exchange} tanpa mengubah \textit{producer}. & VI, IX, X \\
\hline
\textit{Recoverability} & Pendukung & \textit{Queue} durable, \textit{manual acknowledgement}, pod yang dijadwalkan ulang Kubernetes, dan perbaikan salinan shard oleh \textit{repair loop}. Tidak ada \textit{backup} dan retensi. & IV, X, XI \\
\hline
\textit{Reliability} & Pendukung & Pengiriman \textit{at-least-once} dengan kunci deduplikasi di \textit{consumer}. Data bersifat \textit{eventual consistent}. & VII, X \\
\hline
\textit{Testability} & Pendukung & Tiap \textit{service} Go punya \textit{test} sendiri, kedua \textit{mock} menjadi pengganti sumber, dan skrip skenario per problem ada di \texttt{tests/scenarios}. & V \\
\hline
\textit{Deployability} & Pendukung & Satu Dockerfile per \textit{service} dan satu perintah \texttt{deploy/deploy.sh}. \textit{Rolling update} \textit{Aggregator} dengan \texttt{maxUnavailable} 0 (\texttt{docs/DEPLOY.md}). & II, V \\
\hline
\textit{Scalability} & Pendukung & Data dibagi ke 16 shard di node yang dapat bertambah, dan \textit{Aggregator} berjalan 3 sampai 5 replika. \textit{Service} lain tetap satu replika. & XI, XII \\
\hline
\textit{Elasticity} & Pendukung & \textit{Autoscaler} CPU untuk \textit{Aggregator} (\texttt{deploy/aggregator-hpa.yaml}). Node yang bergabung mendapat pod \textit{store} otomatis karena DaemonSet. & XI, XII \\
\hline
\textit{Learnability} & Tidak dikejar & Go dan Svelte dipilih karena mudah dibaca, tetapi tidak ada rancangan khusus untuk anggota baru dan tidak ada pengukuran. & III \\
\hline
\end{longtable}

\textit{Availability} dan \textit{fault tolerance} diutamakan karena spesifikasi meminta satu sumber yang lambat atau mati tidak mengganggu sumber lain. Pada Milestone 1 \textit{scalability} dan \textit{elasticity} sengaja dikorbankan. Di Milestone 2 keduanya hanya dikejar pada \textit{Aggregator} dan \textit{store}, karena problem Milestone 2 menyentuh keduanya. \textit{State} di memori seperti hitungan \textit{rate limit} dan keadaan \textit{circuit breaker} pada \textit{service} lain tetap tidak disinkronkan antar-\textit{instance}, sehingga \textit{service} itu tetap satu replika. Batas ini dicatat pada bagian Batasan yang Disadari di Bab VII, IX, XI, dan XII. \ph{Konfirmasi dengan kelompok bahwa prioritas Scalability dan Elasticity menjadi Pendukung sudah sesuai, lalu selaraskan isi Gambar \ref{fig:karakteristik}}

Gaya arsitektur yang dihasilkan adalah \textit{microservice} dengan penyimpanan di tiap \textit{service} dan, pada Milestone 2, penyimpanan yang dipartisi dan direplikasi serta pemilihan pemimpin berbasis \textit{lease} pada \textit{Aggregator}, ditambah distribusi \textit{event} berpola \textit{publish-subscribe}. Kepercayaan antar-\textit{service} berada di perimeter, sesuai Bab IV.
